%\addcontentsline{toc}{chapter}{Remerciements}

\textit{À tous ceux qui ont contribué, de près ou de loin, à l'aboutissement de ce projet, je dédie ces quelques mots de reconnaissance.}

\vspace{0.8cm}

Au terme de ce projet de fin d'études, il m'est particulièrement agréable d'exprimer ma gratitude envers toutes les personnes dont le soutien, la confiance et l'expertise ont rendu ce travail possible.

\vspace{0.5cm}

Je remercie en premier lieu les membres du jury pour l'honneur qu'ils me font en acceptant d'examiner et d'évaluer ce travail. Leurs retours critiques constitueront une contribution précieuse à ma formation d'ingénieur.

\vspace{0.5cm}

Mes vifs remerciements s'adressent à \textbf{M. Youssef Meddeb}, mon encadrant industriel au sein d'\textbf{Uvey}, pour m'avoir intégré pleinement au sein de l'équipe d'ingénierie dès les premiers jours. Sa pédagogie, sa disponibilité pour répondre à mes questions techniques, et sa confiance dans les choix que j'ai proposés m'ont permis de m'approprier le projet avec un véritable sentiment de responsabilité. Les échanges réguliers en Sprint Review ont été particulièrement formateurs.

\vspace{0.5cm}

Je tiens à remercier \textbf{Mme Meriem Kassar}, mon encadrante académique, pour la qualité de son suivi tout au long de ce projet. Ses remarques constructives, sa disponibilité lors de chaque point d'avancement, et sa capacité à recadrer les priorités ont été déterminantes pour maintenir la cohérence et la rigueur du travail. Je lui suis reconnaissant d'avoir su allier exigence académique et bienveillance pédagogique.

\vspace{0.5cm}

Je remercie chaleureusement \textbf{l'ensemble de l'équipe technique d'Uvey} pour l'accueil bienveillant, les revues de code, les discussions informelles autour des choix d'architecture, et pour m'avoir traité comme un membre à part entière de l'équipe. Ce contexte de travail collaboratif et exigeant a rendu ce projet à la fois stimulant et formateur bien au-delà du cadre académique.

\vspace{0.5cm}



